\chapter{Profil 10 Startup AI Paling Berpengaruh}
\label{chap:startups}

Ekosistem AI didorong oleh inovasi startup yang bergerak cepat. Berikut adalah profil 10 perusahaan yang membentuk lanskap AI modern, dari raksasa riset hingga penyedia infrastruktur.

\section{1. OpenAI}
\textbf{Didirikan:} 2015 (San Francisco, AS) \\
\textbf{Fokus:} Artificial General Intelligence (AGI), LLM, Generative AI. \\
\textbf{Produk Utama:} ChatGPT, GPT-4, DALL-E 3, Sora. \\
\textbf{Pendanaan:} \$13 Miliar+ (Microsoft sebagai investor utama).

\subsection*{Kisah Singkat}
Awalnya didirikan sebagai organisasi non-profit oleh Elon Musk, Sam Altman, dan lainnya untuk memastikan AI bermanfaat bagi umat manusia. Pada 2019, beralih ke struktur "capped-profit" untuk menarik modal besar. OpenAI memicu perlombaan senjata AI global dengan peluncuran ChatGPT pada November 2022, menjadi aplikasi konsumen dengan pertumbuhan tercepat dalam sejarah.

\subsection*{Dampak}
Mempopulerkan Transformer decoder-only (GPT) dan teknik RLHF. Menetapkan standar de facto untuk performa LLM.

\section{2. Anthropic}
\textbf{Didirikan:} 2021 (San Francisco, AS) \\
\textbf{Fokus:} AI Safety, Large Language Models. \\
\textbf{Produk Utama:} Claude (Claude 3 Opus, Sonnet, Haiku). \\
\textbf{Pendanaan:} \$7 Miliar+ (Amazon, Google).

\subsection*{Kisah Singkat}
Didirikan oleh mantan karyawan OpenAI (Dario Amodei, Daniela Amodei) yang keluar karena perbedaan visi mengenai keamanan AI. Anthropic fokus pada "Constitutional AI"---melatih model agar tidak berbahaya dengan memberikan konstitusi etika, bukan hanya feedback manusia (RLHF).

\subsection*{Keunggulan}
Model Claude dikenal memiliki context window yang sangat besar (200k+ token), kemampuan penalaran yang kuat, dan lebih sedikit "halusinasi" berbahaya dibanding kompetitor.

\section{3. Hugging Face}
\textbf{Didirikan:} 2016 (New York, AS / Paris, Prancis) \\
\textbf{Fokus:} Open Source AI, MLOps, Community Hub. \\
\textbf{Produk Utama:} Hugging Face Hub (Model, Dataset, Spaces), Transformers Library. \\
\textbf{Pendanaan:} \$396 Juta (Valuasi \$4.5 Miliar).

\subsection*{Kisah Singkat}
Sering disebut sebagai "GitHub-nya AI". Dimulai sebagai aplikasi chatbot remaja, lalu pivot menjadi platform infrastruktur open source. Hugging Face mendemokratisasi akses ke model SOTA (State-of-the-Art) dengan library `transformers` yang menyatukan PyTorch, TensorFlow, dan JAX.

\subsection*{Dampak}
Menjadi pusat kolaborasi global. Tempat utama untuk hosting model open source seperti Llama, Mistral, dan Stable Diffusion.

\section{4. Midjourney}
\textbf{Didirikan:} 2021 (San Francisco, AS) \\
\textbf{Fokus:} Text-to-Image Generation. \\
\textbf{Produk Utama:} Midjourney (via Discord), Midjourney Alpha (Web). \\
\textbf{Pendanaan:} Bootstrapped (Tanpa modal ventura eksternal, profit sejak awal).

\subsection*{Kisah Singkat}
Didirikan oleh David Holz (co-founder Leap Motion). Unik karena menolak investasi VC besar dan beroperasi dengan tim sangat kecil (~11 orang di awal). Memiliki komunitas seni generatif yang sangat loyal di Discord.

\subsection*{Keunggulan}
Kualitas estetika dan artistik yang dianggap terbaik di kelasnya. Model v6 mereka mampu menghasilkan fotorealisme yang sulit dibedakan dari foto asli.

\section{5. Stability AI}
\textbf{Didirikan:} 2019 (London, Inggris) \\
\textbf{Fokus:} Open Source Generative AI (Image, Audio, Video, Code). \\
\textbf{Produk Utama:} Stable Diffusion, Stable Audio, Stable Video Diffusion. \\
\textbf{Pendanaan:} \$101 Juta+.

\subsection*{Kisah Singkat}
Mengguncang dunia dengan merilis Stable Diffusion secara open source pada 2022, memungkinkan siapa saja menjalankan generator gambar kelas atas di GPU konsumen (lokal). Ini memicu ledakan inovasi komunitas (ControlNet, LoRA) yang tidak mungkin terjadi pada model tertutup.

\subsection*{Tantangan}
Menghadapi tekanan finansial yang berat karena biaya komputasi tinggi untuk melatih model gratis, serta gugatan hak cipta dari seniman. CEO Emad Mostaque mundur pada Maret 2024.

\section{6. Mistral AI}
\textbf{Didirikan:} 2023 (Paris, Prancis) \\
\textbf{Fokus:} Open Weight LLMs, Efisiensi. \\
\textbf{Produk Utama:} Mistral 7B, Mixtral 8x7B (MoE), Mistral Large. \\
\textbf{Pendanaan:} \$640 Juta+ (Valuasi \$6 Miliar dalam <1 tahun).

\subsection*{Kisah Singkat}
Didirikan oleh mantan peneliti DeepMind dan Meta (Llama team). Mereka membuktikan bahwa model kecil (7B parameter) yang dilatih dengan baik bisa mengalahkan model raksasa (Llama 2 13B, GPT-3.5). Mempopulerkan arsitektur Mixture of Experts (MoE) dalam model open weights.

\subsection*{Dampak}
Simbol kedaulatan AI Eropa. Menjadi alternatif utama bagi developer yang butuh model performa tinggi yang bisa di-host sendiri dengan murah.

\section{7. Cohere}
\textbf{Didirikan:} 2019 (Toronto, Kanada) \\
\textbf{Fokus:} Enterprise AI, NLP for Business (RAG, Search, Embedding). \\
\textbf{Produk Utama:} Command R, Embed v3, Rerank. \\
\textbf{Pendanaan:} \$435 Juta+.

\subsection*{Kisah Singkat}
Didirikan oleh Aidan Gomez (salah satu penulis paper "Attention Is All You Need"). Berbeda dengan OpenAI yang mengejar AGI konsumen, Cohere fokus murni pada solusi enterprise: keamanan data, RAG (Retrieval-Augmented Generation), dan model embedding multilingual yang superior.

\subsection*{Keunggulan}
Model "Rerank" mereka adalah standar industri untuk meningkatkan akurasi sistem pencarian/RAG dengan mengurutkan ulang dokumen hasil retrieval.

\section{8. Perplexity AI}
\textbf{Didirikan:} 2022 (San Francisco, AS) \\
\textbf{Fokus:} AI Search Engine. \\
\textbf{Produk Utama:} Perplexity (Search Interface). \\
\textbf{Pendanaan:} \$100 Juta+ (Jeff Bezos, NVIDIA).

\subsection*{Kisah Singkat}
Menantang dominasi Google Search dengan "Answer Engine". Alih-alih memberikan daftar link biru, Perplexity membaca hasil pencarian dan merangkum jawaban langsung dengan sitasi (footnote). Menggunakan model gabungan (GPT-4, Claude 3, Mistral) untuk menghasilkan jawaban terbaik.

\subsection*{Dampak}
Mengubah perilaku pencarian informasi online dari "keyword search" menjadi "conversational search".

\section{9. Databricks (MosaicML)}
\textbf{Didirikan:} 2013 (Databricks), 2020 (MosaicML - Diakuisisi 2023) \\
\textbf{Fokus:} Data Platform, Enterprise LLM Training. \\
\textbf{Produk Utama:} Dolly, MPT (Mosaic Pretrained Transformer), DBRX. \\
\textbf{Status:} MosaicML dibeli Databricks seharga \$1.3 Miliar.

\subsection*{Kisah Singkat}
Databricks (pencipta Apache Spark) mengakuisisi MosaicML untuk memungkinkan perusahaan melatih LLM mereka sendiri dari awal dengan biaya yang jauh lebih murah (resep pelatihan yang dioptimalkan). Merilis model open source DBRX yang sangat kuat.

\subsection*{Filosofi}
"Your Data, Your Model". Mendorong desentralisasi AI agar setiap perusahaan punya model sendiri, bukan bergantung pada API tertutup.

\section{10. ElevenLabs}
\textbf{Didirikan:} 2022 (New York, AS / London, Inggris) \\
\textbf{Fokus:} AI Voice Generation (Text-to-Speech, Dubbing). \\
\textbf{Produk Utama:} Speech Synthesis, Voice Cloning, Dubbing Studio. \\
\textbf{Pendanaan:} \$101 Juta (Valuasi Unicorn).

\subsection*{Kisah Singkat}
Didirikan oleh mantan engineer Google dan Palantir. Mengubah industri sulih suara dengan teknologi yang mampu meniru emosi, intonasi, dan bahkan bahasa pembicara asli dengan sampel suara yang sangat pendek (few-shot).

\subsection*{Dampak}
Memungkinkan konten kreator menjangkau audiens global (dubbing otomatis) dan memberikan suara pada karakter game/buku audio dengan biaya fraksional dibanding studio rekaman tradisional.

\begin{tipbox}
    \textbf{Honorable Mentions:}
    \begin{itemize}
        \item \textbf{Scale AI:} "Pabrik" pelabelan data (RLHF) di balik layar kesuksesan OpenAI dan perusahaan besar lainnya.
        \item \textbf{Character.AI:} Platform chatbot hiburan dengan engagement time tertinggi (rata-rata 2 jam/user/hari).
        \item \textbf{CoreWeave:} Penyedia cloud GPU khusus yang menyediakan infrastruktur untuk Microsoft dan Inflection AI.
    \end{itemize}
\end{tipbox}
